Оценка долгосрочного эффекта онлайн-экспериментов признаётся одной из наиболее
сложных задач как в академических исследованиях, так и в индустриальной
практике. Классические подходы к А/Б-тестированию хорошо приспособлены для
измерения краткосрочных эффектов воздействий, однако с трудом улавливают
отложенные и устойчивые сдвиги в ключевых бизнес-метриках, таких как
пожизненная ценность клиента (customer lifetime value, LTV или CLV)
\cite{crook2009seven,kohavi2020trustworthy}. Ряд работ показывает, что
оптимизация исключительно под краткосрочные исходы рискованна: ранние
выигрыши могут со временем уменьшаться или сменяться потерями вследствие
изменения паттернов вовлечённости пользователей~\cite{sadeghi2021novelty,
kohavi2015focus,kohavi2016pitfalls,dmitriev2016dirty}.

Чтобы преодолеть этот разрыв, в литературе развивается несколько направлений.
Одно из наиболее заметных --- использование \textit{суррогатных} (proxy) метрик,
когда краткосрочные наблюдаемые индикаторы применяются для предсказания
долгосрочных исходов. Так, в исследованиях Google предложены подходы к
построению оптимальных прокси-метрик на основе истории прошлых А/Б-тестов:
показано, что композитные суррогаты заметно улучшают принятие решений по
сравнению с одиночной метрикой~\cite{tripuraneni2024}. В работах LinkedIn
продемонстрированы методы аппроксимации долгосрочных метрик по краткосрочным
ковариатам, позволяющие сокращать длительность экспериментов при сохранении
предсказательной способности~\cite{duan2021online}. Netflix развивает это
направление, обучая прокси-метрики на исторических экспериментах и подчёркивая
важность согласованности направления краткосрочных сигналов и долгосрочных
целей~\cite{dimakopoulou2023evaluating}.

Параллельное направление связано с поведенческой сегментацией клиентской базы.
Классический RFM-анализ (Recency, Frequency, Monetary value) на протяжении
десятилетий используется как базовый инструмент сегментации: он группирует
пользователей по характеристикам транзакционного поведения и даёт интуитивно
понятные и применимые на практике инсайты для бизнес-стратегии
\cite{fader2005rfm,rosset2003modeling}. Дальнейшее развитие этих идей --- модели
жизненного цикла клиента (Pareto/NBD, BG/NBD) и подходы к описанию динамики
поведения, в том числе скрытые марковские модели отношений с
клиентом~\cite{schmittlein1987,schmittlein1987customer,fader2005counting,fader2010,netzer2008mm}.

В настоящей работе предлагается подход, в котором долгосрочные эффекты
А/Б-экспериментов оцениваются через \textit{целеориентированную} поведенческую
сегментацию пользователей. Каждый сегмент определяется условиями над
пользователем и историей его активности в системе; сами условия подбираются
так, чтобы внутри сегмента группировались пользователи с похожими целевыми
характеристиками (например, GMV, удержанием, интенсивностью вовлечённости).
Такие сегменты мы называем \textit{орбиталями} и анализируем не столько
сами их статические характеристики, сколько \textit{траектории} пользователей
между ними. В отличие от статических когорт, такая динамическая сегментация
постоянно обновляет принадлежность пользователей к сегментам, преодолевая
ограничения традиционных подходов, не учитывающих временную
динамику~\cite{REUTTERER200643,zhuzhel2021cohortney,orduz2023cohort}.
По духу метод близок к работам, в которых поведение пользователей описывается
как скрытый марковский процесс~\cite{netzer2008mm,pfeifer2003modeling}, однако
вместо фиксированных скрытых состояний используются непосредственно
орбитали, а наблюдаемым итогом эксперимента становится финальное распределение
пользователей по ним.

В работе формулируются три ключевых результата:
\begin{itemize}
  \item вводится понятие \textit{орбиталей} --- целеориентированной
        поведенческой сегментации, обеспечивающей надёжную и интерпретируемую
        оценку долгосрочных эффектов по коротким экспериментам;
  \item формализуется использование \textit{динамики переходов} между
        орбиталями как механизма причинного вывода, связывающего ближайшие
        поведенческие изменения с будущей ценностью;
  \item метод валидируется на реальных данных А/Б-тестов крупномасштабной
        e-commerce платформы, и демонстрируется, что орбитали улучшают
        качество принимаемых решений, не нарушая темпа быстрых
        экспериментальных циклов.
\end{itemize}

Предлагаемый подход закрывает важный разрыв в современном экспериментировании:
рассогласованность между быстрыми циклами разработки продукта и стратегической
потребностью смотреть на долгосрочные эффекты. Опираясь на наблюдаемые
траектории пользователей, орбитали дают прозрачную и удобную в эксплуатации
альтернативу непрозрачным суррогатным моделям и при этом не требуют
проведения длительных экспериментов.

\begin{comment}
Структура работы соответствует логике решения поставленной задачи. В главе
<<Постановка задачи>> на конкретном производственном примере мотивируется
рабочая интерпретация долгосрочного эффекта и формулируются цель и задачи
исследования. В обзоре литературы рассматриваются основные направления оценки
долгосрочного эффекта онлайн-экспериментов и моделирования поведения
пользователей. Далее вводится математическая модель, лежащая в основе метода
Orbitals: определение ограниченной LTV, орбиталей и структуры переходов
между ними. Затем описывается сам метод --- конструкция орбиталей, вычисление
метрик эффекта и архитектура решения, --- после чего приводятся результаты
экспериментального исследования на реальных А/Б-тестах и сравнительный
анализ с альтернативными подходами. В заключении подводятся итоги работы,
обсуждаются ограничения предложенного подхода и намечаются направления
дальнейших исследований.
\end{comment}